\BourbakiFrontMatter{致读者}

新版

\begin{enumerate}
\def\labelenumi{\arabic{enumi}.}
\item
  《数学原理》丛书从最起点讲起数学并给出完整的证明。因此，原则上它不要求读者具备任何特定的数学知识，只要求对数学推理有一定的熟悉以及一定的抽象思维能力。尽管如此，它主要面向那些至少较好地掌握大学数学头一两年内容的人。
\item
  所选的讲述方法是公理化的，并且通常从一般走向特殊。证明的要求使题材的顺序被严格固定。由此，某些考虑的用处可能要到较后的章节才会显现在读者面前，除非他们已经具有相当广博的数学知识。
\item
  丛书分为若干卷（书），每卷又分为若干章。下面列出在法文版中已经出版（全部或部分）的各卷。当有英译本时，在括号中给出相应的英文标题。在本卷中，引用在可用时指英文版，否则指法文版。
\end{enumerate}

Théorie des ensembles (Theory of Sets) Ref. E (Set Theory)

Algèbre (Algebra \(^{(1)}\) ) --- A (Alg.)

Topologie générale (General Topology) --- TG (Gen.~Top.)

Fonctions d'une variable réelle

(Functions in One Real Variable) --- FVR (FRV)

Espaces vectoriels topologiques

(Topological Vector Spaces)

---

EVT

(Top. Vect. Sp.)

Intégration (Integration)

---

INT

(Int.)

Algèbre commutative

\((Commutative\ Algebra^{(2)})\)

---

AC

(Comm. Alg.)

Variétés différentiables et analytiques

---

VAR

Groupes et algèbres de Lie

(Lie Groups and Lie Algebras)

---

LIE

(Lie)

Théories spectrales

---

TS

Topologie Algébrique

---

TA

在前六卷（按上面提到的顺序）中，正文里的每个陈述都只预设同一章或前面各章中已经讨论过的定义和结果为已知，这些章的先后顺序如下：集合论；Alg.，第 I 至 III 章；Gen.~Top.，第 I 至 III 章；Alg./A，第 IV 章起；Gen.~Top.，第 IV 章起；FRV；Top. Vect. Sp.；Int.。从第七卷起，读者在必要时可在每卷或每章的开头找到对所用到的其他卷或章的精确说明（前六卷总假定已知）。

\begin{enumerate}
\def\labelenumi{\arabic{enumi}.}
\setcounter{enumi}{3}
\item
  然而，某些段落不遵循上述规则。它们被置于两个星号之间：\(*\ldots\) 。在某些情形，这样做只是为了通过收入一些引用读者可能从别处得知的结果的例子来帮助读者理解正文。在另一些情形，我们不仅使用有关各章中预设为已知的结果，还使用丛书中别处证明的结果。这些段落在那些预设包含它们的各章及它们所引用的各章为已知的部分中被自由使用。我们希望读者能够验证其中不存在任何恶性循环。
\item
  某些卷（已出版或在编写中）可以包含结果摘要。这些摘要包含该卷的本质性定义和结果，不含任何证明。
\item
  每章的逻辑骨架由该章的定义、公理和定理组成。这些是此后使用时须牢记的主要部分。较次要的结果，以及那些容易由定理重新推出的结果，冠以``命题''、``引理''、``推论''、``注''等名称。那些初读时可以跳过的部分以小号字体排印。关于某个基本定理的评论偶尔以``评注''（scholium）之名出现。
\end{enumerate}

为避免令人厌倦的重复，有时方便的是引入仅在某一特定的章、节或小节内有效的记号或缩写（例如，在某章中所有环都交换时，``环''一词总可以指``交换环''）。此类约定在它们适用的章、节或小节的开头明确说明。

请注意，与本书前面各章一样，除非另有说明，体不假设是交换的。

\begin{enumerate}
\def\labelenumi{\arabic{enumi}.}
\setcounter{enumi}{6}
\item
  某些段落旨在预先提醒读者警惕严重的错误。这些段落以记号 Z（``危险弯道''）在页边标出。
\item
  习题一方面旨在让读者检验自己是否已很好地消化正文，另一方面让他们熟悉那些在正文中没有位置的结果。最难的习题以 ¶ 标记。
\item
  我们对本丛书所用的术语给予了特别的注意。除似乎有充分理由偏离通行术语之处外，我们都遵循普遍接受的术语。
\item
  我们特别努力做到始终使用严格正确的语言而不牺牲简洁。在正文中，我们尽可能对各种滥用语言或滥用记号之处提请读者注意，没有这些滥用，任何数学文本都有变得迂腐、更谈不上可读的风险。
\item
  由于正文是对一个理论的教条式阐述，它很少包含对文献的引用。书目式引用有时汇集于历史笔记中。这些笔记之后的参考文献一般只包含在对所论理论的演进最为重要的书籍和原始论著。它并不求完备。
\end{enumerate}

至于习题，我们认为没有必要指出其来源，因为它们取自许多不同的来源（原始论文、教科书、习题集）。

\begin{enumerate}
\def\labelenumi{\arabic{enumi}.}
\setcounter{enumi}{11}
\tightlist
\item
  在新版本书的本章中，对定理、公理、定义、注等的引用，是通过依次指出其所在的卷（使用第 3 条中列出的缩写）、章、节、小节和页码而给出的。当所引的卷与本书相同时，则省略对卷的指明。例如，在《代数》一书中，
\end{enumerate}

\textbf{集合论，III，§ 4，No.~2，p.~167，推论 3}

指《集合论》一书第 III 章 § 4 第 2 号第 167 页的推论 3；

\textbf{II，§ 1，No.~11，p.~215，命题 17}

指《代数》一书第 II 章 § 1 第 11 号第 215 页的命题 17。

结果摘要以字母 R 标示；例如 Top. Vect. Sp., R 指《拓扑向量空间》一书的``结果摘要''。

由于某些卷在新版中稍后才出版，对这些卷的引用由依次给出的卷、章、节、小节（所论结果应在之处）组成，例如，

\textbf{Comm. Alg.，III，§ 4，No.~5，命题 6 的推论.}

当引用法文版的一卷时，使用正体大写缩写，并使用法文术语和排印方式；例如，

\textbf{TA, II, § 2, n° 4, p.~158, corollaire de la proposition 1}

指《Topologie Algébrique（代数拓扑）》一书第 II 章 § 2 第 4 号第 158 页的命题 1 的推论。

